\section{Gizmo e o bytecode que o Quarkus gera}
\label{sec:gizmo}

O código que as extensões geram (fábricas de beans, a classe principal,
a reprodução dos recorders) é escrito direto em bytecode com o
\termo{Gizmo}{Biblioteca de geração de bytecode do projeto Quarkus.}. O
README do projeto (\citetitle{doc:gizmo}) o define como \enquote{A
bytecode generation library} e diz que ele é muito usado pelo Quarkus; o
guia de extensões o apresenta como \enquote{a convenient alternative to
ASM, with a higher-level API}. Não há \texttt{.java} intermediário: o
resultado é \texttt{.class}, o mesmo formato que o \texttt{javac} produz e
que você aprendeu a ler no capítulo~\ref{cap:lendo-bytecode}.

\subsection*{Como ver as classes geradas}

Em dev mode e nos testes, as classes geradas e transformadas ficam só na
memória de um class loader, diz o guia \citetitle{doc:quarkus-writing-extensions}. Para
gravá-las em disco ele oferece:

\small
\begin{longtable}{@{}p{5.6cm} p{6.6cm}@{}}
\toprule
Propriedade & O que grava \\
\midrule \endhead
\texttt{quarkus.debug.\allowbreak generated-classes-dir} & Classes geradas, como os
  metadados de beans \\
\texttt{quarkus.debug.\allowbreak transformed-classes-dir} & Classes transformadas,
  por exemplo entidades Panache \\
\texttt{quarkus.debug.\allowbreak generated-sources-dir} & Arquivos ZIG, uma
  representação textual do código gerado \\
\texttt{quarkus.package.jar.\allowbreak decompiler.enabled} & Baixa o descompilador
  Vineflower e grava o resultado no diretório \texttt{decompiled} da
  saída do build; só funciona num build de produção normal (não em dev
  mode nem em testes) e com o empacotamento \emph{fast-jar}, o padrão \\
\bottomrule
\end{longtable}
\normalsize

Com as classes em disco, \texttt{javap -c -p} funciona como em qualquer
outra classe. Esta apostila não reproduz listagens de classes geradas
pelo Quarkus: elas mudam entre versões e não são saída do nosso
\texttt{examples/}. O desafio no fim do capítulo pede que você as gere no
seu projeto.

\subsection*{Por que gerar código em vez de usar reflexão}

O exemplo abaixo é Java puro, compilado com o JDK~25 como o resto da
apostila. Ele mostra as duas formas de um container colocar um
\texttt{Repositorio} dentro de um \texttt{Servico}: por reflexão, num
campo \texttt{private}, ou por uma atribuição comum, num campo
\emph{package-private}, que é o que uma classe gerada no mesmo pacote
consegue fazer.

\lstinputlisting[language=Java]{generated/source/InjecaoSemReflexao.java}

\inputbytecode{InjecaoSemReflexao}

\paragraph{Lendo a listagem.}
\begin{itemize}
\item \texttt{porReflexao} faz três chamadas \texttt{invokevirtual}:
  procura o campo pelo \emph{nome} (a string \texttt{"privado"} no
  constant pool), desliga a checagem de acesso e escreve. O campo
  escrito é um \emph{argumento} de método, não o operando de uma
  instrução. Aqui o argumento é uma constante, e o Native Image consegue
  resolvê-lo no build (seção~\ref{sec:reflexao-nativo}); um container
  que descobre beans varrendo o classpath em runtime, porém, só conhece
  os nomes durante a execução.
\item \texttt{direto} é um único \texttt{putfield} cujo operando aponta
  para o campo no constant pool. Qualquer ferramenta que lê o
  \texttt{.class} (o verificador, o JIT, a análise estática do Native
  Image) sabe qual campo é escrito sem executar nada.
\end{itemize}

\paragraph{A decisão que isso muda.}
É exatamente o argumento do guia \citetitle{doc:quarkus-cdi-reference}:
se a injeção precisa acessar um membro \texttt{private}, o Quarkus tem de
usar reflexão; em executável nativo, todo membro acessado por reflexão
precisa ser registrado. Por isso o guia recomenda não usar membros
\texttt{private} em beans e preferir campos \emph{package-private} ou
injeção pelo construtor. No nativo, o guia aponta o custo: esses registros
de reflexão aumentam o executável (seção~\ref{sec:reflexao-nativo}). Para
campos de injeção \texttt{private}, o padrão
(\texttt{quarkus.arc.transform-private-injected-fields=true}) já reescreve o
campo como \emph{package-private} no build, e a injeção fica livre de
reflexão; construtores, inicializadores, observers, producers, disposers e
métodos de interceptor \texttt{private} continuam exigindo reflexão.

\paragraph{Panache, um caso de transformação.}
Além de gerar classes novas, o Quarkus reescreve classes suas. O guia
\citetitle{doc:quarkus-panache} recomenda campos \texttt{public} nas
entidades e diz que o Panache gera os getters e setters que faltam,
apoiado em \enquote{compile-time bytecode enhancements}. São essas
classes que \texttt{quarkus.debug.transformed-classes-dir} grava.
